%%%%%%%%%%%%%%%%%%%%%%%%%%%%%%%%%%%%%%%%%
%
% CMPT 435
% Lab Zero
%
%%%%%%%%%%%%%%%%%%%%%%%%%%%%%%%%%%%%%%%%%

%%%%%%%%%%%%%%%%%%%%%%%%%%%%%%%%%%%%%%%%%
% Short Sectioned Assignment
% LaTeX Template
% Version 1.0 (5/5/12)
%
% This template has been downloaded from: http://www.LaTeXTemplates.com
% Original author: % Frits Wenneker (http://www.howtotex.com)
% License: CC BY-NC-SA 3.0 (http://creativecommons.org/licenses/by-nc-sa/3.0/)
% Modified by Alan G. Labouseur  - alan@labouseur.com, and Patrick Tyler - Patrick.Tyler1@marist.edu
%
%%%%%%%%%%%%%%%%%%%%%%%%%%%%%%%%%%%%%%%%%

%----------------------------------------------------------------------------------------
%	PACKAGES AND OTHER DOCUMENT CONFIGURATIONS
%----------------------------------------------------------------------------------------

\documentclass[letterpaper, 10pt]{article} 

\usepackage[english]{babel} % English language/hyphenation
\usepackage{graphicx}
\usepackage[lined,linesnumbered,commentsnumbered]{algorithm2e}
\usepackage{listings}
\usepackage{fancyhdr} % Custom headers and footers
\pagestyle{fancyplain} % Makes all pages in the document conform to the custom headers and footers
\usepackage{lastpage}
\usepackage{url}
\usepackage{xcolor}
\usepackage{titlesec}
\usepackage{xurl}
\usepackage{multicol}

% Stolen from https://www.overleaf.com/learn/latex/Code_listing 
\definecolor{codegreen}{rgb}{0,0.6,0}
\definecolor{codegray}{rgb}{0.5,0.5,0.5}
\definecolor{codepurple}{rgb}{0.58,0,0.82}
\definecolor{backcolour}{rgb}{0.95,0.95,0.92}

\lstdefinestyle{mystyle}{
    backgroundcolor=\color{backcolour},   
    commentstyle=\color{codegreen},
    keywordstyle=\color{magenta},
    numberstyle=\tiny\color{codegray},
    stringstyle=\color{codepurple},
    basicstyle=\ttfamily\footnotesize,
    breakatwhitespace=false,         
    breaklines=true,                 
    captionpos=b,                    
    keepspaces=true,                 
    numbers=left,                    
    numbersep=5pt,                  
    showspaces=false,                
    showstringspaces=false,
    showtabs=false,                  
    tabsize=2
}
\lstset{style=mystyle, language=c++}


\fancyhead{} % No page header - if you want one, create it in the same way as the footers below
\fancyfoot[L]{} % Empty left footer
\fancyfoot[R]{}

\renewcommand{\headrulewidth}{0pt} % Remove header underlines
\renewcommand{\footrulewidth}{0pt} % Remove footer underlines
\setlength{\headheight}{13.6pt} % Customize the height of the header

%----------------------------------------------------------------------------------------
%	TITLE SECTION
%----------------------------------------------------------------------------------------

\newcommand{\horrule}[1]{\rule{\linewidth}{#1}} % Create horizontal rule command with 1 argument of height

\title{	
   \normalfont \normalsize 
   \textsc{CMPT 308 - Spring 2024 - Dr. Labouseur} \\[10pt] % Header stuff.
   \horrule{0.5pt} \\[0.25cm] 	% Top horizontal rule
   \huge Lab 7 -- ~Normalization \\     	    % Assignment title
   \horrule{0.5pt} \\[0.25cm] 	% Bottom horizontal rule
}

\author{Patrick Tyler \\ \normalsize Patrick.Tyler1@marist.edu}

\date{\normalsize\today} 	% Today's date.

\begin{document}

\maketitle % Print the title

%----------------------------------------------------------------------------------------
%   CONTENT SECTION
%----------------------------------------------------------------------------------------

% - -- -  - -- -  - -- -  -
\section{Part one}
\subsection{Original table}
\begin{table}[h]
\centering
\begin{tabular}{|c|c|c|c|}
\hline
PackageId & TagNumber & InstallDate & SoftwareCostUSD \\
\hline
AC01 & 32808 & 09-13-2005 & 754.95 \\
\cline{1-4}
DB32 & 32808 & 12-03-2005 & 380.00 \\
     & 37691 & 06-15-2005 & 380.00 \\
\cline{1-4}
DB33 & 57772 & 05-27-2005 & 412.77 \\
\cline{1-4}
WP08 & 32808 & 01-12-2006 & 185.00 \\
     & 37691 & 06-15-2005 & 227.50 \\
     & 57222 & 05-27-2005 & 170.24 \\
\cline{1-4}
WP09 & 59836 & 10-30-2005 & 35.00 \\
     & 77740 & 05-27-2005 & 35.00 \\
\hline
\end{tabular}
\end{table}
\subsection{My professional consulting opinion}
Tycho manufacuting's board needs to fire the CEO 
    as a software company I expected better. 
Asumming this data is stored in a csv, excel, or some
    other raw text table structure to store the spreadsheet,
    the data is difficult to programatically add, update, delete
    additional installs and packages.
If an install needs to be deleted for whatever reason, the
    row with that install would have to have its internal structure
    of some columns edited to remove the correct install 
    without affecting others.
Additionally, nulls may also need to be left to store packages
    without installs.
Accidental deletions of packages from deleting installs could 
    also be a problem.
In my professional opinion, this data is not complex enough
    to warrant a database, but it would help organizing
    the spreadsheet more atomically.
\subsection{Spreasheet in 1NF}
\begin{table}[h]
\centering
\begin{tabular}{|c|c|c|c|}
\hline
PackageId & TagNumber & InstallDate & SoftwareCostUSD \\
\hline
AC01 & 32808 & 09-13-2005 & 754.95 \\
\cline{1-4}
DB32 & 32808 & 12-03-2005 & 380.00 \\
\cline{1-4}
DB32 & 37691 & 06-15-2005 & 380.00 \\
\cline{1-4}
DB33 & 57772 & 05-27-2005 & 412.77 \\
\cline{1-4}
WP08 & 32808 & 01-12-2006 & 185.00 \\
\cline{1-4}
WP08 & 37691 & 06-15-2005 & 227.50 \\
\cline{1-4}
WP08 & 57222 & 05-27-2005 & 170.24 \\
\cline{1-4}
WP09 & 59836 & 10-30-2005 & 35.00 \\
\cline{1-4}
WP09 & 77740 & 05-27-2005 & 35.00 \\
\hline
\end{tabular}
\end{table}
\subsection{Primary key analaysis}
The best primary key that can be made with these columns
    is an composite key of the PackageID, TagNumber, 
    and InstallDate.
This has the side affect of dissallowing mutliple installs
    that have the same TagNumber and PackageID on the same day.
\section{Part two adding Package Name and Computer Model}
\subsection{New Table}
\begin{table}[h]
\centering
\begin{tabular}{|c|c|c|c|c|c|}
\hline
PackageId & PackageName & TagNumber & ComputerModel & InstallDate & SoftwareCostUSD \\
\hline
AC01 & Zork & 32808 & Lenova & 09-13-2005 & 754.95 \\
\cline{1-6}
DB32 & Roket Leauge & 32808 & Lenova & 12-03-2005 & 380.00 \\
\cline{1-6}
DB32 & Rocket Leauge & 37691 & Apple & 06-15-2005 & 380.00 \\
\cline{1-6}
DB33 & Chocolately & 57772 & IBM & 05-27-2005 & 412.77 \\
\cline{1-6}
WP08 & Neovim & 32808 & Lenova & 01-12-2006 & 185.00 \\
\cline{1-6}
WP08 & Neovim & 37691 & Apple & 06-15-2005 & 227.50 \\
\cline{1-6}
WP08 & Neovim & 57222 & IBM & 05-27-2005 & 170.24 \\
\cline{1-6}
WP09 & Balatro & 59836 & Dell & 10-30-2005 & 35.00 \\
\cline{1-6}
WP09 & Balatro & 77740 & MSI & 05-27-2005 & 35.00 \\
\hline
\end{tabular}
\end{table}
\subsection{Functional Dependencies}
The functional dependencies are sets of columns where some
    column(s) are depedent on their value(s) from other column(s).
In this table, PackageName is dependent on the PackageId and
    ComputerModel is depdent on TagNumber which makes these two
    functional dependenices.
\subsection{Non 3NF-ness}
This table is not in third normal form because a part,
    importantly not the whole, of the primary key
    functionally determines another column.
Both of the aformentioned functional dependencies create
    this problem as, for example, with just the PackageID
    you already know what the PackageName will be.
Practically, this creates duplication within the table
    because you are storing identifying information of packages
    twice.
You could just have another table which maps the PackageId
    to its PackageName.
This way you would only need to write the PackageName once
    in the spreadsheet.
\end{document}
